\documentclass[11pt]{article}
\usepackage[margin=1in]{geometry}
\usepackage{amsmath,amssymb,amsthm}
\usepackage{enumitem}
\usepackage{booktabs}

\newtheorem{theorem}{Theorem}
\newtheorem{proposition}[theorem]{Proposition}
\newtheorem{corollary}[theorem]{Corollary}
\newtheorem{remark}[theorem]{Remark}

\title{d5\_247: explicit two-swap splice closes \(G1\)}
\author{}
\date{2026-03-22}

\begin{document}
\maketitle

\section{Purpose}
The remaining graph-side object \(G1\) asked for a full colorwise package
containing the already closed color-\(4\) branch.
Earlier candidate searches showed that several narrow ans\"atze fail:
naive cyclic copies fail, mode-switch families fail, and same-color coarse-class
splices on support of size \(3\) or \(4\) fail.
The point of the present note is that \(G1\) nevertheless has a simple explicit
solution once one stops insisting that the non-color-\(4\) repairs depend only
on their \emph{own} color-relative class.
The correct package is a two-defect-set transposition of the valid baseline.

\section{Baseline package}
Write \(X_m=(\mathbb Z/m\mathbb Z)^5\), and for
\(x=(x_0,x_1,x_2,x_3,x_4)\in X_m\) define
\[
\sigma(x):=x_0+x_1+x_2+x_3+x_4 \pmod m.
\]
For each color \(c\in\{0,1,2,3,4\}\), let \(B_c\colon X_m\to X_m\) be the
explicit baseline package from d5\_245:
\[
B_c(x)=x+e_{d_c(x)},
\qquad
d_c(x)=
\begin{cases}
c+1 \pmod 5, & \sigma(x)=0,\\
c-1 \pmod 5, & \sigma(x)=1,\\
c, & \sigma(x)\notin\{0,1\}.
\end{cases}
\]
This package is pointwise outgoing-exhaustive and colorwise permutation for every
\(m>2\).

\section{The color-\(4\) defect sets}
The surfaced \texttt{mixed\_008} color-\(4\) rule differs from \(B_4\) on
exactly two source families.

\medskip
\noindent\textbf{First defect set.}
Define
\[
D_1:=\{x\in X_m:\ \sigma(x)=2,\ x_0=m-1\}.
\]
This is exactly the color-\(4\) class \(L2^{(1)}\), i.e.\ the layer-\(2\) old-bit
\(1\) family for the surfaced mixed branch.
On \(D_1\), the mixed rule uses direction \(1\) instead of the baseline
direction \(4\).

\medskip
\noindent\textbf{Second defect set.}
For \(m=3\), set \(D_2:=\varnothing\).
For \(m\ge 5\), define
\[
D_2:=\bigl\{x\in X_m:\ \sigma(x)=3,\ 
(x_0\neq m-1\ \wedge\ x_1+x_3=2)\ \vee\
(x_0=m-1\ \wedge\ x_1+x_3\neq 2)\bigr\}.
\]
This is exactly the union of the surfaced color-\(4\) classes
\[
L3^{(0,p=1)}\ \cup\ L3^{(1,p=0)}.
\]
On \(D_2\), the mixed rule uses direction \(2\) instead of the baseline
direction \(4\).

\begin{proposition}[mixed color-\(4\) rule as a two-set swap]
\label{prop:mixed-color4}
For every odd \(m>2\), the surfaced \texttt{mixed\_008} color-\(4\) rule is
\[
\mathrm{SelStar}_4(x)=
\begin{cases}
x+e_1, & x\in D_1,\\
x+e_2, & x\in D_2,\\
B_4(x), & x\notin D_1\cup D_2.
\end{cases}
\]
\end{proposition}

\begin{proof}
At layer \(0\) and layer \(1\), the surfaced mixed rule has the same anchors as
the baseline, namely anchor \(1\) on \(L0\) and anchor \(4\) on \(L1\), so the
absolute directions for color \(4\) are again \(0\) and \(3\), matching \(B_4\).
At layers \(4+\), the anchor is again \(0\), so the absolute direction is \(4\),
matching \(B_4\).

At layer \(2\), the mixed table is \(0/2\): old-bit \(0\) keeps anchor \(0\)
(direction \(4\)), while old-bit \(1\) takes anchor \(2\) (direction \(1\)).
For color \(4\), the old bit is \(x_0=m-1\), so the exceptional set is exactly
\(D_1\).

At layer \(3\), the mixed tables are \(0/3\) on predecessor flag \(p=0\) and
\(3/0\) on predecessor flag \(p=1\).  Therefore the absolute color-\(4\)
direction changes from \(4\) to \(2\) exactly on the two classes
\(L3^{(0,p=1)}\) and \(L3^{(1,p=0)}\).  For color \(4\), the old bit is
\(x_0=m-1\), and the simple predecessor flag \(\texttt{pred\_sig1\_wu2}\) is
equivalent to \(x_1+x_3=2\), because the predecessor in relative direction \(1\)
changes the \(q=x_0\) coordinate only and leaves \(w=x_1\) and \(u=x_3\)
unchanged.
Thus the union of those two exceptional classes is exactly \(D_2\).
\end{proof}

\section{The explicit splice package}
Define maps \(g_0,\dots,g_4\colon X_m\to X_m\) by
\[
g_0:=B_0,\qquad g_3:=B_3,
\]
and
\[
g_1(x)=
\begin{cases}
x+e_4, & x\in D_1,\\
B_1(x), & x\notin D_1,
\end{cases}
\qquad
g_2(x)=
\begin{cases}
x+e_4, & x\in D_2,\\
B_2(x), & x\notin D_2,
\end{cases}
\]
while
\[
g_4(x)=
\begin{cases}
x+e_1, & x\in D_1,\\
x+e_2, & x\in D_2,\\
B_4(x), & x\notin D_1\cup D_2.
\end{cases}
\]
So relative to the baseline, color \(1\) and color \(4\) are swapped on \(D_1\),
and color \(2\) and color \(4\) are swapped on \(D_2\).

\section{Image-preserving identities}
The construction is powered by two translation invariances.

\begin{proposition}[defect-set invariance]
\label{prop:defect-invariance}
For every odd \(m>2\):
\begin{enumerate}[label=\textup{(\arabic*)},leftmargin=2.8em]
\item \(D_1\) is invariant under translation by \(e_1-e_4\);
\item \(D_2\) is invariant under translation by \(e_2-e_4\).
\end{enumerate}
\end{proposition}

\begin{proof}
For \(D_1\), the defining conditions are \(\sigma(x)=2\) and \(x_0=m-1\).
Translation by \(e_1-e_4\) preserves the total sum and does not touch \(x_0\),
so it preserves \(D_1\).

For \(D_2\), the defining conditions (when \(m\ge 5\)) are
\(\sigma(x)=3\), a condition on \(x_0\), and a condition on
\(x_1+x_3\).
Translation by \(e_2-e_4\) preserves the total sum, leaves \(x_0\) fixed, and
leaves \(x_1+x_3\) fixed.
Hence it preserves \(D_2\).
For \(m=3\), \(D_2=\varnothing\), so the statement is trivial.
\end{proof}

\begin{corollary}[image preservation]
\label{cor:image-preserving}
For every odd \(m>2\),
\[
B_1(D_1)=B_4(D_1),
\qquad
B_2(D_2)=B_4(D_2).
\]
\end{corollary}

\begin{proof}
If \(x\in D_1\), then by Proposition~\ref{prop:defect-invariance} also
\(x+e_1-e_4\in D_1\), and
\[
B_4(x+e_1-e_4)=x+e_1=B_1(x).
\]
So \(B_1(D_1)\subseteq B_4(D_1)\).  The reverse inclusion is identical with
\(e_4-e_1\).

Likewise, if \(x\in D_2\), then \(x+e_2-e_4\in D_2\) and
\[
B_4(x+e_2-e_4)=x+e_2=B_2(x),
\]
which gives \(B_2(D_2)=B_4(D_2)\).
\end{proof}

\section{Closure of \(G1\)}
\begin{theorem}[explicit two-swap splice]
\label{thm:G1-closed}
For every odd \(m>2\), the maps \(g_0,\dots,g_4\) defined above satisfy:
\begin{enumerate}[label=\textup{(\arabic*)},leftmargin=2.8em]
\item outgoing exhaustiveness holds pointwise;
\item each \(g_c\) is a permutation of \(X_m\);
\item \(g_4\) is exactly the surfaced closed color-\(4\) rule
      \(\mathrm{SelStar}_4=\texttt{mixed\_008}\).
\end{enumerate}
Hence the graph-side frontier theorem \(G1\) is closed.
\end{theorem}

\begin{proof}
By Proposition~\ref{prop:mixed-color4}, item \textup{(3)} is immediate.

For outgoing exhaustiveness, outside \(D_1\cup D_2\) the package equals the
baseline package \(B\), so the property is unchanged.
On \(D_1\), the only modification is the transposition of the color-\(1\) and
color-\(4\) directions: \(B_1\) and \(B_4\) are exchanged, so the five outgoing
directions remain \(\{0,1,2,3,4\}\).
On \(D_2\), the same is true for the transposition of colors \(2\) and \(4\).
Since \(D_1\cap D_2=\varnothing\) (they lie in different layers), item
\textup{(1)} follows.

For the permutation statement, \(g_0=B_0\) and \(g_3=B_3\) are already
permutations.
For \(g_1\), the defect set is \(D_1\), the modified defect map is \(B_4\), and
Corollary~\ref{cor:image-preserving} gives
\[
g_1(D_1)=B_4(D_1)=B_1(D_1).
\]
Because \(B_4\vert_{D_1}\) is injective, the finite-defect splice principle
implies that \(g_1\) is a permutation.
Exactly the same argument, with \(D_2\), shows that \(g_2\) is a permutation.

For \(g_4\), the defect set is \(D_1\cup D_2\), and
\[
g_4(D_1\cup D_2)=B_1(D_1)\cup B_2(D_2)=B_4(D_1)\cup B_4(D_2)=B_4(D_1\cup D_2),
\]
again by Corollary~\ref{cor:image-preserving}.
Because \(D_1\) and \(D_2\) are disjoint and \(B_4\) is injective, the two image
sets \(B_4(D_1)\) and \(B_4(D_2)\) are disjoint, so \(g_4\vert_{D_1\cup D_2}\)
is injective.
Thus \(g_4\) is a permutation as well.
This proves item \textup{(2)}.
\end{proof}

\begin{remark}[why the earlier searches still mattered]
The earlier no-go results were not wasted:
they ruled out several \emph{too narrow} ans\"atze.
In particular, same-color coarse-class repairs fail.
The successful splice above is different: the correction for color \(1\) is
driven by the color-\(4\) defect set \(D_1\), and the correction for color \(2\)
is driven by the color-\(4\) defect set \(D_2\).
So the winning package is still explicit and local, but it is not a same-color
classwise splice.
\end{remark}

\section{Checked computational certificate}
The accompanying checker \texttt{d5\_247\_G1\_explicit\_splice\_checker.py}
verifies on checked odd moduli
\[
m=3,5,7,9,11,13
\]
that:
\begin{enumerate}[label=\textup{(\arabic*)},leftmargin=2.8em]
\item the surfaced color-\(4\) defect sets agree with the coordinate formulas
      for \(D_1\) and \(D_2\);
\item the image-preserving identities
      \(B_1(D_1)=B_4(D_1)\) and \(B_2(D_2)=B_4(D_2)\) hold exactly;
\item the explicit package is outgoing-exhaustive and colorwise Latin; and
\item the resulting color-\(4\) map agrees pointwise with the surfaced
      \texttt{mixed\_008} rule.
\end{enumerate}
This computation is only a certificate; the proof above is symbolic and uniform
for all odd \(m>2\).

\end{document}
